\section{Conclusion}
\label{sec:conclusion}

We conducted a systematic evaluation of text reproduction fidelity in three GPT-4.1-family models across five categories of stimuli and approximately 500 test cases. Our central finding is that the landscape of LLM ``tongue twisters'' has shifted dramatically: classic \glitchtokens are completely resolved, and models faithfully preserve misspellings without auto-correcting. However, new failure modes have emerged at extreme text lengths. Most notably, \gptfull paradoxically truncates clean text at $\geq$8k characters while perfectly reproducing 25k characters of misspelled text---a content-dependent asymmetry that suggests output filtering rather than a capability limitation. All models also fail on long repetitive text, with \gptmini producing output 5.15$\times$ the input length in the most severe case.

These findings have direct implications for any application that relies on LLMs for faithful text reproduction, from code copying to data transcription. Future work should test whether the clean-versus-misspelled asymmetry generalizes across model families, investigate the mechanism behind content-dependent truncation, and explore whether chain-of-thought or chunked reproduction strategies can overcome length-based limitations.
